\paragraph{The improved arithmetic density and final delta bounds.}
The established uniform divisor now has exponent profile
\[
 g_r(x)+\mathbf1_{\mathcal S}(\{x\})
       +\mathbf1_{\mathcal T}(\{x\}),\qquad r=3,4,
\]
where $g_r$ is the baseline profile, $\mathcal S$ is the fourteen-band
first-order set, and
\[
 \mathcal T=(1/10,2/19)\cup(11/12,15/16).
\]
The second-order lemma supplies exactly the additional factor used in the
last indicator; it is added only once. The nine certified second-order
chambers lie inside the first-order set. Their valuations are $c+3$, while
the baseline and first-order valuations are $c+1$ and $c+2$, respectively.

All prime restrictions leave the limiting density unchanged. For
$p^2\le30n$, the logarithm of their product is $O(\sqrt n)=o(n)$.
Fixed excluded primes contribute $O(1)$, and primes dividing $n$ contribute
at most $\log n$. The profile exponents are bounded. Every nonzero
$\{n/p\}$ for $p>29$ has reduced denominator $p$, so it cannot equal an
endpoint or omitted Farey-$29$ wall. For any fixed number of periods, the
prime number theorem evaluates their density. The remaining periods are
bounded by $O(1/K)$, since they use only primes $p\le n/K$; this justifies
passing to the full positive series.

The exact finite sums and elementary infinite-tail enclosures give
\[
 \gamma\ge
 \frac{30458547243805493774844991}{10^{24}}
 >30.45854724380549.
\]
In particular we may use the shorter rational bound $\gamma>30.45854724$.
Combining it with the accepted analytic rates
$a<65.451536565866$ and $b<-47.356002681589$, the integer height and
linear-form decay exponents satisfy
\[
\begin{array}{c|cc|c}
 &\sigma&\tau&\liminf_{N\to\infty}\delta_N\\\hline
 B=0,\ A\ne0
 &91.992989325866&20.814549921589&>0.22626234\\
 B\ne0
 &110.992989325866&1.814549921589&>0.01634832.
\end{array}
\]
Indeed $\sigma=a+19r-\gamma$ and $\tau=-b-19r+\gamma$, with $r=3$ for
the first row and $r=4$ for the second. Exact rational cross-multiplication
verifies $\tau/\sigma$ exceeds the displayed decimal in each case.
The index shift $n=N-1$ does not change any exponential rate.

These conclusions hold for every fixed nonzero integer pair $(A,B)$;
fixed rational pairs follow by clearing a fixed denominator. Combining
the two numerator readouts cannot worsen their common error upper bound,
and reducing the resulting rational approximant cannot worsen its delta.
The arithmetic improvement is therefore valid for the stated initial
conditions. The accepted analytic estimates remain explicit inputs.
No irrationality-measure upper bound is asserted here.
